\renewcommand{\theequation}{C\thechapter.\arabic{equation}}
\label{sec:appendixC}

The proposed spatially resolved model will use Moltres to simulate circulation of a fissile sample in molten salt.
The spatially resolved model requires turbulent modeling to capture vortices and other thermal-hydraulic effects that influence the residence time of fissile material in different regions.
In contrast, the 0D scaled and 0D flow models approximate these effects by varying residence times and applying sensitivity analyses to assess the impact of prolonged residence in the ex-core or in-core regions on the \gls{DNP} group parameters.

The proposed transient simulations using Moltres could also include the turbulence effects.
Turbulence would not affect the \gls{DNP} group parameters, but it would affect the concentration distribution for each group.
Since only the \gls{MSRE} is studied in these transients, turbulent modeling may not be necessary to capture the effects of interest.
The in-core region of the \gls{MSRE} consists of flow through channels, which would not have large residence time differences from turbulence.
Turbulence would cause mixing in the upper and lower plena, so the Moltres simulation without turbulence will assume perfect mixing in those regions.

In the proposed transient simulations, neglecting turbulence may affect residence time in the upper and lower plena, slightly increasing the ex-core residence time relative to the in-core residence time.
Because Moltres can model turbulence using the Spalart-Allmaras turbulence model, it is straightforward to include in the transient simulations, though it does introduce more variables \cite{spalart_one-equation_1992, park_advancements_2025}.
The purpose of the transient simulations is to demonstrate the difference between simulations using improved \gls{DNP} group parameters compared to existing group parameters.
Modeling turbulence would not significantly change this difference, but would improve the simulation accuracy.

For these reasons, modeling turbulence in the Moltres transient simulations will be treated as an optional extension.
If sufficient time is available, turbulent modeling will be included.
If a small amount of additional time is available, a sensitivity analysis of the difference between transients with and without turbulence will be conducted.
If insufficient time is available, this analysis will be moved to the future work section.